\section{Affine schemes}
\label{section-affine-schemes}

\noindent
Let $R$ be a ring. Consider the topological space $\Spec(R)$
associated to $R$, see Algebra, Section \ref{algebra-section-spectrum-ring}.
We will endow this space with a sheaf of rings $\mathcal{O}_{\Spec(R)}$
and the resulting pair $(\Spec(R), \mathcal{O}_{\Spec(R)})$
will be an affine scheme.

\medskip\noindent
Recall that $\Spec(R)$ has a basis of open sets $D(f)$,
$f \in R$ which we call standard opens, see Algebra,
Definition \ref{algebra-definition-Zariski-topology}.
In addition, the intersection of two standard opens is another:
$D(f) \cap D(g) = D(fg)$, $f, g\in R$.

\begin{lemma}
\label{lemma-standard-open}
Let $R$ be a ring. Let $f \in R$.
\begin{enumerate}
\item If $g\in R$ and $D(g) \subset D(f)$, then
\begin{enumerate}
\item $f$ is invertible in $R_g$,
\item $g^e = af$ for some $e \geq 1$ and $a \in R$,
\item there is a canonical ring map $R_f \to R_g$, and
\item there is a canonical $R_f$-module map
$M_f \to M_g$ for any $R$-module $M$.
\end{enumerate}
\item Any open covering of $D(f)$ can be refined to a finite
open covering of the form $D(f) = \bigcup_{i = 1}^n D(g_i)$.
\item If $g_1, \ldots, g_n \in R$, then $D(f) \subset \bigcup D(g_i)$
if and only if $g_1, \ldots, g_n$ generate the unit ideal in $R_f$.
\end{enumerate}
\end{lemma}

\begin{proof}
Recall that $D(g) = \Spec(R_g)$ (see
Algebra, Lemma \ref{algebra-lemma-standard-open}).
Thus (a) holds because $f$
maps to an element of $R_g$ which is not
contained in any prime ideal, and hence invertible,
see Algebra, Lemma \ref{algebra-lemma-Zariski-topology}.
Write the inverse of $f$ in $R_g$ as $a/g^d$.
This means $g^d - af$ is annihilated by a power of $g$, whence (b).
For (c), the map $R_f \to R_g$ exists by (a) from the universal property
of localization, or we can define it by mapping $b/f^n$
to $a^nb/g^{ne}$. The equality $M_f = M \otimes_R R_f$
can be used to obtain the map on modules, or
we can define $M_f \to M_g$ by mapping
$x/f^n$ to $a^nx/g^{ne}$.

\medskip\noindent
Recall that $D(f)$ is quasi-compact, see
Algebra, Lemma \ref{algebra-lemma-qc-open}.
Hence the second statement follows directly
from the fact that the standard opens form
a basis for the topology.

\medskip\noindent
The third statement follows directly from
Algebra, Lemma \ref{algebra-lemma-Zariski-topology}.
\end{proof}

\noindent
In Sheaves, Section \ref{sheaves-section-bases} we defined
the notion of a sheaf on a basis, and we showed that it is
essentially equivalent to the notion of a sheaf on the space,
see Sheaves, Lemmas \ref{sheaves-lemma-extend-off-basis} and
\ref{sheaves-lemma-extend-off-basis-structures}. Moreover,
we showed in
Sheaves, Lemma \ref{sheaves-lemma-cofinal-systems-coverings-standard-case}
that it is sufficient to check the sheaf
condition on a cofinal system of open coverings for each
standard open. By the lemma above it suffices to check
on the finite coverings by standard opens.

\begin{definition}
\label{definition-standard-covering}
Let $R$ be a ring.
\begin{enumerate}
\item A {\it standard open covering} of $\Spec(R)$
is a covering $\Spec(R) = \bigcup_{i = 1}^n D(f_i)$,
where $f_1, \ldots, f_n \in R$.
\item Suppose that $D(f) \subset \Spec(R)$ is a standard
open. A {\it standard open covering} of $D(f)$
is a covering $D(f) = \bigcup_{i = 1}^n D(g_i)$,
where $g_1, \ldots, g_n \in R$.
\end{enumerate}
\end{definition}

\noindent
Let $R$ be a ring. Let $M$ be an $R$-module. We will define
a presheaf $\widetilde M$ on the basis of standard opens.
Suppose that $U \subset \Spec(R)$ is a standard open.
If $f, g \in R$ are such that $D(f) = D(g)$, then
by Lemma \ref{lemma-standard-open} above there are canonical
maps $M_f \to M_g$ and $M_g \to M_f$ which are mutually inverse.
Hence we may choose any $f$ such that $U = D(f)$
and define
$$
\widetilde M(U) = M_f.
$$
Note that if $D(g) \subset D(f)$, then by
Lemma \ref{lemma-standard-open} above we have
a canonical map
$$
\widetilde M(D(f)) = M_f \longrightarrow M_g = \widetilde M(D(g)).
$$
Clearly, this defines a presheaf of abelian groups on the basis
of standard opens. If $M = R$, then $\widetilde R$ is a presheaf
of rings on the basis of standard opens.

\medskip\noindent
Let us compute the stalk of $\widetilde M$ at a point $x \in \Spec(R)$.
Suppose that $x$ corresponds to the prime $\mathfrak p \subset R$.
By definition of the stalk we see that
$$
\widetilde M_x = \colim_{f\in R, f\not\in \mathfrak p} M_f
$$
Here the set $\{f \in R, f \not \in \mathfrak p\}$ is preordered by
the rule $f \geq f' \Leftrightarrow D(f) \subset D(f')$.
If $f_1, f_2 \in R \setminus \mathfrak p$, then we have
$f_1f_2 \geq f_1$ in this ordering. Hence by
Algebra, Lemma \ref{algebra-lemma-localization-colimit}
we conclude that
$$
\widetilde M_x = M_{\mathfrak p}.
$$

\medskip\noindent
Next, we check the sheaf condition for the standard open coverings.
If $D(f) = \bigcup_{i = 1}^n D(g_i)$, then the sheaf condition
for this covering is equivalent with the exactness of the
sequence
$$
0 \to M_f \to \bigoplus M_{g_i} \to \bigoplus M_{g_ig_j}.
$$
Note that $D(g_i) = D(fg_i)$, and hence we can rewrite this
sequence as the sequence
$$
0 \to M_f \to \bigoplus M_{fg_i} \to \bigoplus M_{fg_ig_j}.
$$
In addition, by Lemma \ref{lemma-standard-open} above
we see that $g_1, \ldots, g_n$ generate the unit ideal
in $R_f$. Thus we may apply
Algebra, Lemma \ref{algebra-lemma-cover-module}
to the module $M_f$ over $R_f$ and the elements $g_1, \ldots, g_n$.
We conclude that the sequence is exact. By the remarks
made above, we see that $\widetilde M$ is a sheaf
on the basis of standard opens.

\medskip\noindent
Thus we conclude from the material in
Sheaves, Section \ref{sheaves-section-bases}
that there exists a
unique sheaf of rings $\mathcal{O}_{\Spec(R)}$
which agrees with $\widetilde R$ on the standard opens.
Note that by our computation of stalks above, the
stalks of this sheaf of rings are all local rings.

\medskip\noindent
Similarly, for any $R$-module $M$ there exists
a unique sheaf of $\mathcal{O}_{\Spec(R)}$-modules
$\mathcal{F}$ which agrees with $\widetilde M$ on the
standard opens, see
Sheaves, Lemma \ref{sheaves-lemma-extend-off-basis-module}.

\begin{definition}
\label{definition-structure-sheaf}
Let $R$ be a ring.
\begin{enumerate}
\item The {\it structure sheaf $\mathcal{O}_{\Spec(R)}$ of the
spectrum of $R$} is the unique sheaf of rings $\mathcal{O}_{\Spec(R)}$
which agrees with $\widetilde R$ on the basis of standard opens.
\item The locally ringed space
$(\Spec(R), \mathcal{O}_{\Spec(R)})$ is called
the {\it spectrum} of $R$ and denoted $\Spec(R)$.
\item The sheaf of $\mathcal{O}_{\Spec(R)}$-modules
extending $\widetilde M$ to all opens of $\Spec(R)$
is called the sheaf of $\mathcal{O}_{\Spec(R)}$-modules
associated to $M$. This sheaf is denoted $\widetilde M$ as
well.
\end{enumerate}
\end{definition}

\noindent
We summarize the results obtained so far.

\begin{lemma}
\label{lemma-spec-sheaves}
Let $R$ be a ring. Let $M$ be an $R$-module. Let $\widetilde M$
be the sheaf of $\mathcal{O}_{\Spec(R)}$-modules
associated to $M$.
\begin{enumerate}
\item We have $\Gamma(\Spec(R), \mathcal{O}_{\Spec(R)}) = R$.
\item We have $\Gamma(\Spec(R), \widetilde M) = M$ as an $R$-module.
\item For every $f \in R$ we have
$\Gamma(D(f), \mathcal{O}_{\Spec(R)}) = R_f$.
\item For every $f\in R$ we have $\Gamma(D(f), \widetilde M) = M_f$
as an $R_f$-module.
\item Whenever $D(g) \subset D(f)$ the restriction mappings
on $\mathcal{O}_{\Spec(R)}$ and $\widetilde M$
are the maps
$R_f \to R_g$ and $M_f \to M_g$ from Lemma
\ref{lemma-standard-open}.
\item Let $\mathfrak p$ be a prime of $R$, and let $x \in \Spec(R)$
be the corresponding point. We have
$\mathcal{O}_{\Spec(R), x} = R_{\mathfrak p}$.
\item Let $\mathfrak p$ be a prime of $R$, and let $x \in \Spec(R)$
be the corresponding point. We have $\widetilde M_x = M_{\mathfrak p}$
as an $R_{\mathfrak p}$-module.
\end{enumerate}
Moreover, all these identifications are functorial in the $R$
module $M$. In particular, the functor $M \mapsto \widetilde M$
is an exact functor from the category of $R$-modules
to the category of $\mathcal{O}_{\Spec(R)}$-modules.
\end{lemma}

\begin{proof}
Assertions (1) - (7) are clear from the discussion above.
The exactness of the functor $M \mapsto \widetilde M$
follows from the fact that the functor $M \mapsto M_{\mathfrak p}$
is exact and the fact that exactness of short exact sequences
may be checked on stalks, see
Modules, Lemma \ref{modules-lemma-abelian}.
\end{proof}

\begin{definition}
\label{definition-affine-scheme}
An {\it affine scheme} is a locally ringed space isomorphic
as a locally ringed space to $\Spec(R)$ for some ring $R$.
A {\it morphism of affine schemes} is a morphism in the category
of locally ringed spaces.
\end{definition}

\noindent
It turns out that affine schemes play a special role
among all locally ringed spaces, which is what the next
section is about.


















